\documentclass[10pt]{article}

\usepackage[a4paper,margin=0.62in]{geometry}
\usepackage[T1]{fontenc}
\usepackage[utf8]{inputenc}
\usepackage{enumitem}
\usepackage{tabularx}
\usepackage{xcolor}
\usepackage{microtype}
\usepackage[
  unicode=true,
  colorlinks=true,
  linkcolor=black,
  citecolor=black
]{hyperref}

\definecolor{ink}{HTML}{1F2937}
\definecolor{muted}{HTML}{5F6B7A}
\definecolor{rule}{HTML}{D9DEE8}
\definecolor{accent}{HTML}{5B3F8C}
\hypersetup{urlcolor=accent}

\setlength{\parindent}{0pt}
\setlength{\parskip}{0pt}
\setlength{\emergencystretch}{2em}
\renewcommand{\familydefault}{\sfdefault}

\newcommand{\cvsection}[1]{%
  \vspace{0.75em}
  {\large\bfseries\color{accent} #1}\par
  \vspace{0.2em}
  {\color{rule}\hrule height 0.6pt}
  \vspace{0.45em}
}

\newcommand{\entry}[3]{%
  \textbf{#1}\hfill{\color{muted}#2}\\
  {\color{muted}#3}\par
}

\begin{document}
\pagestyle{empty}
\color{ink}

\begin{tabularx}{\linewidth}{@{}Xr@{}}
  {\LARGE\bfseries Junhuai Xu} &
  {\large\bfseries\color{accent}Curriculum Vitae} \\
  Ph.D. Candidate, Department of Physics, Tsinghua University &
  {\small\color{muted}Updated July 2026} \\
  \textbf{Ph.D. Advisor}: \href{https://inspirehep.net/authors/1062622}{Zhigang Xiao} &
  \href{mailto:xjh22@mails.tsinghua.edu.cn}{xjh22@mails.tsinghua.edu.cn} \\
  Beijing, China &
  \href{https://orcid.org/0000-0001-7529-8786}{ORCID} \quad
  \href{https://inspirehep.net/authors/1866896}{INSPIRE-HEP} \quad
  \href{https://scholar.google.com/citations?user=XK-LVeIAAAAJ\&hl=en}{Google Scholar}
\end{tabularx}

\vspace{0.55em}
{\color{accent}\hrule height 1.0pt}

\cvsection{Profile}
Ph.D. candidate in experimental nuclear and particle physics, working across detector instrumentation,
scientific computing, and physics-informed reconstruction. I develop analysis workflows that connect detector
calibration, response simulation, event reconstruction, inverse problems, statistical inference, and
machine-learning methods for nuclear and neutrino experiments.

\cvsection{Research Experience}
\entry{Nucleon Short-Range Correlations}{Oct. 2023 -- Present}{CSHINE Experiment, RIBLL-1 (HIRFL), Lanzhou, China}
\begin{itemize}[leftmargin=1.25em,itemsep=1pt,topsep=2pt]
  \item Extracted the SRC-induced high-momentum-tail fraction in ${}^{124}$Sn, \(R_{\mathrm{HMT}}=(20\pm3)\%\), by comparing precision hard-$\gamma$ spectra with detector-filtered IBUU-MDI transport calculations.
  \item Developed calibration and detector-response models for the CsI(Tl) $\gamma$-ray hodoscope used in bremsstrahlung measurements from low-energy heavy-ion collisions.
  \item Implemented Richardson--Lucy spectrum unfolding and quantified reconstruction/systematic uncertainties through pseudo-experiments, closure tests, and consistency checks.
\end{itemize}

\entry{Correlation Function Imaging}{Feb. 2024 -- Present}{Department of Physics, Tsinghua University}
\begin{itemize}[leftmargin=1.25em,itemsep=1pt,topsep=2pt]
  \item Developed a Richardson--Lucy imaging framework for femtoscopic two-particle correlations, reconstructing freeze-out source distributions without assuming a Gaussian source shape.
  \item Reconstructed identical non-Gaussian source functions for $pp$ and $\bar{p}\bar{p}$ correlations in Au+Au collisions at $\sqrt{s_{NN}}=200$ GeV, providing coordinate-space evidence for matter--antimatter symmetry at freeze-out.
  \item Extended the imaging strategy toward three-dimensional source reconstruction with reproducible fit, resampling, and validation workflows for collaboration analyses.
\end{itemize}

\entry{Deep-Sea Neutrino Detection}{Jun. 2024 -- Present}{TRIDENT Collaboration}
\begin{itemize}[leftmargin=1.25em,itemsep=1pt,topsep=2pt]
  \item Built an atmospheric-neutrino simulation-to-reconstruction chain combining HONDA fluxes, FLUKA interactions, and Geant4 detector transport for a compact TRIDENT array.
  \item Developed graph-neural-network reconstruction of event class, direction, and energy from simulated detector responses, benchmarking performance for oscillation analyses.
  \item Evaluated atmospheric-neutrino oscillation-parameter sensitivity using Poisson-sampled pseudo-experiments and spectral comparisons across parameter hypotheses.
\end{itemize}

\cvsection{Experimental Experience}
\begin{itemize}[leftmargin=1.25em,itemsep=2pt,topsep=2pt]
  \item \textbf{CSHINE Experiment} (Lanzhou, China, 2024): participated in beam-time data taking and high-energy $\gamma$-ray measurements with the CSHINE detector system.
  \item \textbf{SLEGS Beam Test at SSRF} (Shanghai, China, Jan. 2025): performed high-energy $\gamma$-ray response calibration of CsI(Tl) crystals using quasi-monochromatic photon beams and Geant4 simulations.
  \item \textbf{S$\pi$RIT Experiment} (RIKEN, Japan, 2024): supported detector operation and electronics control during SAMURAI beam time, contributing to stable data taking and run coordination.
\end{itemize}

\cvsection{Technical Skills}
\begin{tabularx}{\linewidth}{@{}lX@{}}
\textbf{Physics analysis} & Detector calibration, detector-response simulation, event reconstruction, spectrum unfolding, systematic uncertainty evaluation.\\
\textbf{Inference} & Inverse problems, Richardson--Lucy deconvolution/imaging, likelihood-based inference, pseudo-experiments, closure tests, resampling.\\
\textbf{Machine learning} & Graph neural networks, PyTorch, physics-informed reconstruction, simulation-to-inference workflows.\\
\textbf{Computing} & Python, C/C++, ROOT, Geant4, FLUKA-based workflows, Linux, Git, HPC/GPU-accelerated computing, LaTeX.\\
\end{tabularx}

\newpage
\cvsection{Education}
\entry{Tsinghua University}{Aug. 2022 -- Jun. 2027 (Expected)}{Ph.D. Candidate in Physics}
Research areas: short-range correlations, heavy-ion collisions, femtoscopy, neutrino simulation, and machine-learning-based reconstruction.\\[0.25em]
\entry{South China Normal University}{Aug. 2018 -- Jun. 2022}{B.S. in Optoelectronic Information Science and Engineering}
GPA: 4.06/5.00; Rank: 1/137. Undergraduate thesis in DIS/SIDIS phenomenology: \textit{A Study of the Nuclear Medium Effect of Hadron Fragmentation Function}.

\cvsection{Selected Publications}
{\small A complete publication list is available separately on my personal website.}
\begin{itemize}[leftmargin=1.25em,itemsep=2pt,topsep=3pt]
  \item \textbf{Xu, J.} \textit{et al.}, ``Precise measurement of short-range correlations in nuclei from bremsstrahlung gamma-ray emission in low-energy heavy-ion collisions,'' \href{https://doi.org/10.1103/jw1p-36pb}{\textbf{\textit{Phys. Rev. Res.}} \textbf{7}, 043174 (2025)}.
  \item \textbf{Xu, J.} \textit{et al.}, ``Experimental study of bremsstrahlung $\gamma$-ray emission and short-range correlations in ${}^{124}$Sn+${}^{124}$Sn collisions at 25 MeV/nucleon,'' \href{https://doi.org/10.1103/dhz2-nl56}{\textbf{\textit{Phys. Rev. C}} \textbf{113}, 044613 (2026)}.
  \item \textbf{Xu, J.} \textit{et al.}, ``Reconstruction of Bremsstrahlung $\gamma$-rays spectrum in heavy ion reactions with Richardson-Lucy algorithm,'' \href{https://doi.org/10.1016/j.physletb.2024.139009}{\textbf{\textit{Phys. Lett. B}} \textbf{857}, 139009 (2024)}.
  \item \textbf{Xu, J.} \textit{et al.}, ``Imaging Freeze-Out Sources and Extracting Strong Interaction Parameters in Relativistic Heavy-Ion Collisions,'' \href{https://doi.org/10.1088/0256-307X/42/3/031401}{\textbf{\textit{Chin. Phys. Lett.}} \textbf{42}, 031401 (2025)}.
  \item \textbf{Xu, J.} \textit{et al.}, ``Linear response of CsI(Tl) crystal to energetic photons below 20 MeV,'' \href{https://doi.org/10.1016/j.nima.2025.170787}{\textbf{\textit{Nucl. Instrum. Meth. A}} \textbf{1080}, 170787 (2025)}.
\end{itemize}

\cvsection{Awards}
\begin{itemize}[leftmargin=1.25em,itemsep=2pt,topsep=2pt]
  \item First-class Comprehensive Scholarship, Tsinghua University (2025).
  \item Second-class Comprehensive Scholarship, Tsinghua University (2024).
  \item Outstanding Teaching Assistant, Tsinghua University (2024).
  \item Third Prize, National Undergraduate Mathematics Competition Final Round, Non-Math Major (2021).
  \item First Prize, China Undergraduate Mathematical Contest in Modeling (CUMCM, 2020).
\end{itemize}

\end{document}
